\subsection{\texorpdfstring{向量空间 \(C^{n}\)}{向量空间 C 的 n 次幂}}

由于拓扑空间 \(\mathbf{C}\) 可与 \(\mathbf{R}^2\) 等同，拓扑积 \(\mathbf{C}^n\) （ \(n\) 个与 \(\mathbf{C}\) 全同的空间之积）作为拓扑空间可与 \(\mathbf{R}^{2n}\) 等同；同样，\(\mathbf{C}^n\) 的拓扑群结构（即 \(n\) 个因子的加法（拓扑）群结构之积）可与加法群 \(\mathbf{R}^{2n}\) 的结构等同。但由于 \(\mathbf{C}\) 是一个域，我们可以在 \(\mathbf{C}^n\) 上定义 \(n\) 维向量空间（在 \(\mathbf{C}\) 上）的结构：积 \(az\) （复数 \(a\) 与点 \(z = (z_i)\) （属于 \(\mathbf{C}^n\) ）的积）是点 \((az_i)\) ；这个向量空间结构应仔细地与 \(2n\) 维向量空间（ \(\mathbf{R}\) 上的、定义在 \(\mathbf{R}^{2n}\) 上的）结构区别开来（第六章，§ 1，第 3 号）。我们将把记号 \(\mathbf{C}^n\) 保留给 \(n\) 个与 \(\mathbf{C}\) 全同的空间之积所成的拓扑空间，它还带有刚才定义的 \(\mathbf{C}\) 上的向量空间结构；\(\mathbf{C}^n\) 称为 \(n\) 维复数空间。注意，映射 \((t, z) \to tz\) 在 \(\mathbf{C} \times \mathbf{C}^n\) 上是连续的。

\(C^{n}\) 到 \(C^{m}\) 中的仿射线性映射也是 \(R^{2n}\) 到 \(R^{2m}\) 中的仿射线性映射，但反之不真。

例如，映射 \(z \rightarrow \overline{z}\) 是向量空间 \(R^{2}\) 到其自身上的线性映射，但不是向量空间 C 到其自身上的线性映射。

因此，\(C^{n}\) 到 \(C^{m}\) 中的每个仿射线性映射都一致连续；特别地，\(C^{n}\) 到其自身上的每个仿射线性映射都是同胚。

每个 p 维 \((p \leqslant n)\) 的仿射线性簇（在向量空间 \(C^{n}\) 中）也是向量空间 \(R^{2n}\) 中的一个 2p 维仿射线性簇；这里同样，反之不真。为避免一切混淆，\(C^{n}\) 中的 p 维（仿射）线性簇将称为 p 维复线性簇（\(R^{2n}\) 中的线性簇在为避免误解而需要时称为实线性簇）。特别地，一维（相应地二维）复线性簇称为复直线（相应地复平面），而 n - r 维复线性簇称为复超平面。

把实数空间 \(R^{n}\) 看作嵌入在复数空间 \(C^{n}\) 之中，这常常是方便的，其作法是把 \(R^{n}\) 与 \(C^{n}\) 的一个子集（由关系式 \(\mathfrak{J}(z_{k}) = o\) （ \(1 \leqslant k \leqslant n\) ）所定义的）等同起来。\(C^{n}\) 的拓扑群结构在这个子集上诱导的拓扑群结构与 \(R^{n}\) 的拓扑群结构一致。

注意，\(R^{n}\) （这样嵌入在 \(C^{n}\) 中的）不是 \(C^{n}\) 中的复线性簇。

\(R^{n}\) 中 p 个向量组成的一个系，若在 R 上是自由的，则在 C 上也是自由的。\(R^{n}\) 中每个 p 维实线性簇 V 生成一个 p 维复线性簇 \(V'\) （在 \(C^{n}\) 中），使得 V 是 \(V'\) 在 \(R^{n}\) 上的迹；若 V 由 n-p 个线性方程 \(f_{k}(x)=a_{k}\) 组成的一个方程组所定义，其中诸 \(f_{k}\) 是 \(R^{n}\) 上的线性型（具有实系数且在 R 上线性无关），而诸 \(a_{k}\) 是实数，则同样的方程组定义 \(V'\) ，只是现在 x 的坐标取复数值。

反之，若一个 \(p\) 维复线性簇与 \(\mathbf{R}^n\) 有非空的交，则这个交是一个实线性簇，但其维数可以 \(< p\) 。

